% ==========================================================================
% Formatvorlage f�r Diplomarbeiten der HS Karlsruhe und Furtwangen
% Erstellt von Stefan Jassen, Stefan R�mmele-Werner und Manuel Sautermeister
% ==========================================================================

%==================================================================================
% Themeauswahl
% htwg    - HTWG Konstanz
%==================================================================================
\newcommand{\layoutpres}{htwg}

% ==========================================================================
% Meta-Informationen
% Informationen �ber das Dokument, wie z.B. Titel, Autor, Matrikelnr. etc
% ==========================================================================
% Informationen ------------------------------------------------------------
% 	Definition von globalen Parametern, die im gesamten Dokument verwendet
% 	werden k�nnen (z.B auf dem Deckblatt etc.).
% --------------------------------------------------------------------------
\newcommand{\titel}{Angewandte Systemprogrammierung}
\newcommand{\art}{AVR OS}
\newcommand{\Fakultaet}{Name der aktuellen Fakult�t}
\newcommand{\autor}{Andreas Reinhardt}
\newcommand{\autors}{Matthias Weis}
\newcommand{\autorsmall}{Andreas Reinhardt, Matthias Weis}
\newcommand{\Studiengang}{Name des aktuellen Studiengangs}
\newcommand{\matrikelnr}{123456}
\newcommand{\betreuer}{Daniel Urbanietz}
\newcommand{\referent}{Prof.\,Dr. Erstkorrektor}
\newcommand{\korreferent}{Prof.\,Dr. Zweitkorrektor}
\newcommand{\bearbeitungszeit}{Wintersemester 15/16}
\newcommand{\company}{\mbox{XY}}

% Eigene Befehle
\newcommand{\matlab}[1]{\lstinline[basicstyle=\ttfamily\normalsize]|#1|}
\newcommand{\tief}[1]{$_\text{#1}$}			% Tiefstellen
\newcommand{\hoch}[1]{$^\text{#1}$}			% Hochstellen


% Befehle F�r Abk�rzungs- und Symbolverzeichnis\gls*{symb:Lambda}
\newcommand{\symb}[1]{\glslink*{#1}{\gls*{#1}}}	% Gibt Symbol aus
\newcommand{\abk}[1]{\glslink*{#1}{#1}}	% Gibt Abk�rzung aus
\newcommand{\addabk}[2]{\newacronym{#1}{#1}{#2}}	% F�gt Abk�rzung hinzu (nur Pr�ambel)

\newcommand{\addsymb}[3] 	% F�gt Symbol hinzu (nur Pr�ambel)
{
	\newglossaryentry{#1}
	{
		name=#2,
		description={#3},
		sort=symbolphi, type=symbolslist
	}
}

% Befehle zur semantischen Auszeichnung von Text
\newcommand{\Fachbegriff}[1]{\textit{#1}}

% weitere Abk�rzungen
\newcommand{\ua}{\mbox{u.\,a.\ }}
\newcommand{\zb}{\mbox{z.\,B.\ }}
\newcommand{\bs}{$\backslash$}



% ==========================================================================
% Dokumentenkopf
% ==========================================================================
\documentclass[
	%12pt,								% Schriftgr��e kann hier ge�ndert werden
	%DIV=calc,						% Neuberechnung des Satzspiegels		
	ngerman,							% f�r Umlaute, Silbentrennung etc.
	titlepage,						% es wird eine Titelseite verwendet
	a4paper,							% Papierformat
	%headings=normal,			% Schriftgr��e in �berschriften
	numbers=noenddot,			% Nummerierung der �berschrift erh�lt keinen Endpunkt
	listof=totoc,					% Verzeichnisse im Inhaltsverzeichnis auff�hren
	bibliography=totoc,		% Literaturverzeichnis im Inhaltsverzeichnis auff�hren
	index=totoc,					% Index im Inhaltsverzeichnis auff�hren
	parskip=half,					% Absatzabstand (half/full)
	oneside,							% Druck: Einseitig: oneside; Doppelseitig: twoside
												% (passt Rand und Kopf-/Fu�zeile an)
	chapterprefix=true		% "Kapitel 1" statt "1. Kapitel"
	]{scrreprt}
	
% ==========================================================================
% Ben�tigte Packages
% ==========================================================================
% ==========================================================================
% Anpassung Seitenlayout
% ==========================================================================
\usepackage[
	automark,			% Kapitelangaben in Kopfzeile automatisch erstellen
]{scrpage2}

% ==========================================================================
% Anpassung an deutsche Sprache
% ==========================================================================
\usepackage{ngerman}
\usepackage[ngerman]{babel}

% ==========================================================================
% KOMAScriptanpassung (u.a.: Redefine von addtolist)
% ==========================================================================
\usepackage{scrhack}

% ==========================================================================
% f�r compactitem
% ==========================================================================
\usepackage{paralist} 

% ==========================================================================
% Einstellungen f�r �berschriften
% ==========================================================================
\usepackage[bf,big]{titlesec}	% �berschriften schlanker und gr��er %rm,big,compact
\titleformat{\chapter}[display]				% Anpassen des Schriftzugs "Kapitel x"
  {\normalfont\Huge}
  {\chaptertitlename\ \thechapter}{12pt}{\Huge}
  \titlespacing*{\chapter}{0cm}{-\topskip}{30pt}[0pt]
% ==========================================================================
% Einstellung Schriften und Umlaute
% ==========================================================================
\usepackage[ansinew]{inputenc}
\usepackage[T1]{fontenc}
\usepackage{ae} % "sch�neres" �
\usepackage{eurosym}
\usepackage{lmodern}

% ==========================================================================
% SI-Einheiten
% ==========================================================================
\usepackage[locale = DE, binary-units]{siunitx}
\DeclareSIUnit[number-unit-product = {}]\dmips{DMIPS}
\DeclareSIUnit[number-unit-product = {}]\baud{Baud}

\usepackage{icomma}
\usepackage{nicefrac}
% ==========================================================================
% Einbinden von Grafiken
% ==========================================================================
\usepackage[dvips,final]{graphicx}
\graphicspath{{Bilder/}} % Dort liegen die Bilder des Dokuments

% ==========================================================================
% Befehle aus AMSTeX f�r mathematische Symbole z.B. \boldsymbol \mathbb
% ==========================================================================
\usepackage{amsmath,amsfonts}
\usepackage{amsmath, amsthm, amssymb}
\usepackage{mathtools}

% ==========================================================================
% Einfache Definition der Zeilenabst�nde und Seitenr�nder etc.
% ==========================================================================
\usepackage{setspace}
\usepackage{geometry}

% ==========================================================================
% Zum Umflie�en von Bildern
% ==========================================================================
\usepackage{floatflt}
\usepackage{float}    

% ==========================================================================
% Quellcode-Ausgabe
% ==========================================================================
\usepackage{listings}
\usepackage{xcolor} 

% ==========================================================================
% Lange URLs umbrechen etc.
% ==========================================================================
\usepackage[hyphens]{url}

% ==========================================================================
% Einstellungen Bibtex
% ==========================================================================
\usepackage[numbers]{natbib}
\bibliographystyle{alphadin}		%DIN-Stil des Literaturverzeichnisses

% ==========================================================================
% Zum fortlaufenden Durchnummerieren der Fu�noten
% ==========================================================================
\usepackage{chngcntr}

% ==========================================================================
% f�r lange Tabellen
% ==========================================================================
\usepackage{longtable}
\usepackage{array}
\usepackage{ragged2e}
\usepackage{lscape}

% ==========================================================================
% Spaltendefinition rechtsb�ndig mit definierter Breite
% ==========================================================================
\newcolumntype{w}[1]{>{\raggedleft\hspace{0pt}}p{#1}}

% ==========================================================================
% Erweiterung f�r Querverweise
% ==========================================================================
\usepackage[ngerman]{varioref}

% ==========================================================================
% Erweiterung f�r figures
% ==========================================================================
\usepackage[aboveskip=1em,belowskip=0em]{caption}
\usepackage{subfigure}

% ==========================================================================
% Optischer Randausgleich
% ==========================================================================
\usepackage[tracking=true]{microtype}
\DeclareMicrotypeSet*[tracking]{my}{font = */*/*/sc/*}
\SetTracking{ encoding = *, shape = sc }{ 45 }

% ==========================================================================
% Auswahl der Schriftart
% ==========================================================================
% Arial mit euler f�r Mathe:
%\usepackage[scaled=0.92]{uarial}			% Paket f�r Schrift Arial 
%\renewcommand\familydefault{\sfdefault} % Umstellen auf serifenlose Schrift
%\usepackage{eulervm}

% Times mit passendem Mathe-Font
\renewcommand{\rmdefault}{ptm}
\usepackage{times}
\usepackage{txfonts}
% Kapitel im Inhaltsverzeichnis auch in Times darestellen
\setkomafont{chapterentry}{\rmfamily\bfseries\normalsize}	

% ==========================================================================
% griechische Buchstaben
% ==========================================================================
\usepackage{upgreek}	

% ==========================================================================
% Mehrspaltiger Text
% ==========================================================================
\usepackage{multicol}
\columnseprule 1pt	% Linie swischen Spalten
\columnsep 15mm 	% Abstand zwischen Spaletn

% ==========================================================================
% ToDo-List
% ==========================================================================
\usepackage[
%disable,   %hier wird die ToDo List deaktiviert!!!
ngerman,
textsize=footnotesize,
colorinlistoftodos,
figwidth=0.8\textwidth
]{todonotes}

% ==========================================================================
% Draft Markierung
% ==========================================================================
\usepackage[
final, 			%Aktivierung durch draft; Deaktivierung durch final
%draft,
allpages,
color=gray,
mark=\mbox{vorl�ufiger Entwurf zur Durchsicht},
angle=90,
scale=0.08,
xcoord=-96,
ycoord=-35
]{draftmark}

% ==========================================================================
% Sonstiges
% ==========================================================================
\usepackage{pdflscape}
\usepackage{wallpaper}
\usepackage{ifthen}												% Bedingungen

\setlength{\marginparwidth}{2cm}
\reversemarginpar

% ==========================================================================
% Glossar, Abk�rzungsverzeichnis, Symbolverzeichnis
% ==========================================================================
\usepackage[
nonumberlist, %keine Seitenzahlen anzeigen
acronym,      %ein Abk�rzungsverzeichnis erstellen
toc]      %im Inhaltsverzeichnis auf section-Ebene erscheinen
{glossaries}

% Eigene Definition Abk�rzungsverzeichnis:
\newglossarystyle{myacro}
{
	\renewenvironment{theglossary}{}{}
	\renewcommand*{\glossaryheader}{}
	\renewcommand*{\glsgroupheading}[1]{}
	\renewcommand*{\glsgroupskip}{}
	\renewcommand*{\glossaryentryfield}[5]
	{
		\begin{minipage}[t]{0.2\textwidth}
			\glstarget{##1}{##2}~\dotfill
		\end{minipage}
		\begin{minipage}[t]{0.7\textwidth}
			~##3
		\end{minipage}
		\vskip -0.16cm  % Abstand zwischen den Zeilen
	}
	\renewcommand*{\glossarysubentryfield}[6]{%
	\glossaryentryfield{##2}{##3}{##4}{##5}{##6}}%
}

% Eigene Definition Symbolverzeichnis:
\newglossarystyle{mysombol}
{
	\renewenvironment{theglossary}
	{\begin{multicols}{2}\noindent}{\end{multicols}}
	\renewcommand*{\glossaryheader}{}
	\renewcommand*{\glsgroupheading}[1]{}
	\renewcommand*{\glsgroupskip}{}
	\renewcommand*{\glossaryentryfield}[3]
	{\noindent
		\begin{minipage}[t]{0.94\columnwidth}
			\begin{minipage}[t]{0.2\columnwidth}
				\glstarget{##1}{##2}
			\end{minipage}
			\begin{minipage}[t]{0.78\columnwidth}
				##3
			\end{minipage}
		\end{minipage}
		\vskip 0.01em   % Abstand zwischen den Zeilen
	}
	\renewcommand*{\glossarysubentryfield}[3]{%
	\glossaryentryfield{##2}{##3}}%
}
%Ein eigenes Symbolverzeichnis erstellen
\newglossary[slg]{symbolslist}{syi}{syg}{Symbolverzeichnis}
%Den Punkt am Ende jeder Beschreibung deaktivieren
\renewcommand*{\glspostdescription}{}

% ==========================================================================
% PDF-Optionen (Hyperref) (immer zuletzt)
% ==========================================================================
\usepackage[
%colorlinks=false,
bookmarks,
bookmarksopen=true,
pdftitle={\art},
pdfauthor={\autor},
pdfcreator={\autor},
pdfsubject={\titel},
pdfkeywords={\titel},
plainpages=false,% zur korrekten Erstellung der Bookmarks
pdfpagelabels=true,% zur korrekten Erstellung der Bookmarks
hypertexnames=false,% zur korrekten Erstellung der Bookmarks
breaklinks=true
%linktocpage % Seitenzahlen anstatt Text im Inhaltsverzeichnis verlinken
]{hyperref}

\usepackage{bookmark}

% ==========================================================================
% Seiteneinstellungen
% ==========================================================================
% ==========================================================================
% Zeilenabstand
% ==========================================================================
\onehalfspacing

% ==========================================================================
% Farbdefinitionen
% ==========================================================================
\definecolor{grau}{RGB}{240,240,240} %Beispielwert f�r Grau

% ==========================================================================
% Seitenr�nder
% ==========================================================================
\geometry{headsep=0.6cm}				% Abstand Text Oberkante zu Kopfzeile Unterkante
\geometry{footskip=1.1cm}				% Abstand Text Unterkante zu Fu�zeile Oberkante
\geometry{top=3.2cm}						% Abstand Text zu oberem Papierrand
\geometry{bottom=3.7cm}					% Abstand Text zu unterem Papierrand
\geometry{left=3.45cm}					% Abstand Text zu linkem Papierrand
\geometry{right=2.7cm}					% Abstand Text zu rechtem Papierrand

\setlength{\headheight}{1.1\baselineskip}

% ==========================================================================
% Einstellungen Kopf- und Fu�zeile
% ==========================================================================
\pagestyle{scrheadings}
\renewcommand*{\chapterpagestyle}{scrheadings} % Kopf-/Fu�zeile  auf Kapitelanfang
\renewcommand{\headfont}{\normalfont} % Schriftform der Kopfzeile
% Kopfzeile
\ihead[]{\upshape\headmark}
\chead[]{}
\ohead[]{}
\setheadsepline{0.5pt}
% Fu�zeile
\ifthenelse{\equal{\layoutpres}{htwg}}
{
	\refoot[]{\raisebox{-1.0cm}{\includegraphics[height=0.9cm]{Bilder/HTWG_Logo}}}
}{}
\cefoot[]{\raisebox{-0.1cm}{\footnotesize{\upshape\art~\autorsmall}}}

\cofoot[]{\raisebox{-0.1cm}{\footnotesize{\upshape\art~\autorsmall}}}

\rofoot[]{\upshape\pagemark}

\ifthenelse{\equal{\layoutpres}{htwg}}
{
	\lofoot[]{\raisebox{-1.0cm}{\includegraphics[height=0.9cm]{Bilder/HTWG_Logo}}}
}{}
\setfootsepline{0.5pt}

% ==========================================================================
% erzeugt ein wenig mehr Platz hinter einem Punkt
% ==========================================================================
\frenchspacing 

% ==========================================================================
% Nummerierungstiefe ausw�hlen
% ==========================================================================
\setcounter{secnumdepth}{3} % Bis 1.1.1.1
\setcounter{tocdepth}{4} % Bis 1.1.1.1

% ==========================================================================
% Schusterjungen und Hurenkinder vermeiden
% ==========================================================================
\clubpenalty = 10000
\widowpenalty = 10000 
\displaywidowpenalty = 10000

% ==========================================================================
% Definition von Farben f�r Quellcodeausgabe
% ==========================================================================
\definecolor{hellgelb}{rgb}{1,1,0.9}
\definecolor{quellcode_comment}{rgb}{0.2,0.6,0.2}
\definecolor{quellcode_string}{rgb}{0.6,0,0.8}
\definecolor{quellcode_background}{rgb}{0.98,0.98,0.98}
\lstset{numbers=left,
				backgroundcolor=\color{hellgelb},
				breaklines=true,
				frame=single,
				captionpos=b,
 				language=C,
 				basicstyle=\ttfamily\color{black}\scriptsize,
 				keywordstyle=\color{blue},
 				commentstyle=\color{quellcode_comment},
 				stringstyle=\color{quellcode_string}}
 				
\renewcommand{\lstlistingname}{Code} 				

% ==========================================================================
% Fu�noten fortlaufend durchnummerieren
% ==========================================================================
\counterwithout{footnote}{chapter}

% ==========================================================================
% Definitionen des GLossars
% ==========================================================================
%Verwendete Abkürzungen

\addabk{ASIC}{\textbf{A}pplication \textbf{S}pecific \textbf{I}ntegrated \textbf{C}ircuit}
\addabk{SIL}{\textbf{S}afety \textbf{I}ntegrity \textbf{L}evel}
\addabk{UART}{\textbf{U}niversal \textbf{A}synchronous \textbf{R}eceiver/\textbf{T}ransmitter}
\addabk{SPI}{\textbf{S}erial \textbf{P}eriphal \textbf{I}nterface}
\addabk{MISO}{\textbf{M}aster \textbf{I}n \textbf{S}lave \textbf{O}ut}
\addabk{MOSI}{\textbf{M}aster \textbf{O}ut \textbf{S}lave \textbf{I}n}
\addabk{ADC}{\textbf{A}nalog \textbf{D}igital \textbf{C}onverter}
\addabk{DAC}{\textbf{D}igital \textbf{A}nalog \textbf{C}onverter}
\addabk{DSP}{\textbf{D}igitaler \textbf{S}ignal\textbf{p}rozessor}
\addabk{CRC}{\textbf{C}yclic \textbf{R}edundancy \textbf{C}heck}
\addabk{CPU}{\textbf{C}entral \textbf{P}rocessing \textbf{U}nit}
\addabk{DMA}{\textbf{D}irect \textbf{M}emory \textbf{A}ccess}
\addabk{DTC}{\textbf{D}ata \textbf{T}ransfer \textbf{C}ontroller}
\addabk{MIPS}{\textbf{M}illion \textbf{I}nstructions \textbf{P}er \textbf{S}econd}
\addabk{DMIPS}{\textbf{D}hrystone \textbf{MIPS}}
\addabk{RAM}{\textbf{R}andom \textbf{A}ccess \textbf{M}emory}
\addabk{EMV}{\textbf{E}lectro\textbf{m}agnetische \textbf{V}ertr�glichkeit}
\addabk{PAP}{\textbf{P}rogramm\textbf{a}blauf\textbf{p}lan}
\addabk{SCI}{\textbf{S}erial \textbf{C}ommunication \textbf{I}nterface}
\addabk{USART}{\textbf{U}niversial \textbf{S}ynchronous/\textbf{A}synchronous \textbf{R}eceiver/\textbf{T}ransmitter}
\addabk{HART}{\textbf{H}ighway \textbf{A}ddressable \textbf{R}emote \textbf{T}ransducer}
%Symbole
\addsymb{Rv}{\ensuremath{R_\text{v}}}{Vorwiderstand}
\addsymb{f}{\ensuremath{f}}{Frequenz}

\makeglossaries				% Glossar erstellen

\begin{document}

% ==========================================================================
% Definitionen f�r Silbentrennung (Fehlerhafte Trennung von LaTex)
% ==========================================================================
%fuer ggf. falsche Silbentrennungen
\hyphenation{Trenn-bar-es}
\hyphenation{Mess-un-ge-nau-ig-keit}
\hyphenation{mehr-er-er}
\hyphenation{von-ei-nan-der}


% ==========================================================================
% Deckblatt (Auswahl aus zwei verschiedenen)
% ==========================================================================
\pagenumbering{alph}		% Deckblatt erh�llt Seitennummer a, damit Seite 1 nicht doppelt belegt ist 

%\newgeometry{left=1.8cm,right=0.9cm,bottom=1.3cm,top=1.5cm}
\thispagestyle{plain}
\ThisTileWallPaper{\paperwidth}{\paperheight}{Cover/DeckblattAusPDF.pdf}
\begin{titlepage}	
	\hfill
\end{titlepage}
\restoregeometry
 %Deckblatt aus PDF "DeckblattAusPDF.pdf" in Ordner Cover
%\cleardoublepage

\newgeometry{left=1.8cm,right=0.9cm,bottom=1.3cm,top=1.5cm}
\thispagestyle{plain}
\ifthenelse{\equal{\layoutpres}{htwg}}
{
	\ThisTileWallPaper{\paperwidth}{\paperheight}{Bilder/Cover_HTWG.pdf}
}{}
\begin{titlepage}	
		\sffamily
			\begin{minipage}[t!]{0.35\textwidth}
			\begin{flushleft}
				\ifthenelse{\equal{\layoutpres}{htwg}}
				{
					\includegraphics[width=0.75\textwidth]{Bilder/HTWG_Logo}
				}{}
			\end{flushleft}
			\end{minipage}
			\hfill
			\begin{minipage}[t!]{0.5\textwidth}
			\begin{flushright}
			
			\end{flushright}
			\end{minipage}
			%\vfill
			\vskip 1cm
			\hskip 8.5cm
		\begin{minipage}[t!]{0.57\textwidth}
			\begin{flushleft}
				\hfill
				\vskip 2cm
				{\fontsize{45}{10}\textsc{\art}}
				\vskip 2.7cm
				{\Huge\titel\\}
			\end{flushleft}
		\end{minipage}
			\vfill

			\begin{minipage}[!b]{0.38\textwidth}
			\hfill
			\end{minipage}
			\quad
			{\normalsize
			\begin{minipage}[!b]{0.57\textwidth}
				\begin{minipage}[t]{0.38\textwidth}
						Name:\\
						\\
						Betreuer:\\
						Bearbeitungszeit:\\
				\end{minipage}
				\begin{minipage}[t]{0.57\textwidth}
					\autor\\
					\autors\\
					\betreuer\\
					\bearbeitungszeit
				\end{minipage}
			\end{minipage}
			}
%		\end{addmargin}
	\ClearWallPaper
\end{titlepage}
\restoregeometry
 %normales Deckblatt
\cleardoublepage

% ==========================================================================
% Inhaltsverzeichnis
% ==========================================================================
\pagenumbering{Roman}		% Seitennummer in r�mischen Ziffern
\chapter*{Abstract}
Im Rahmen der Lehrveranstaltung angewandte Systemprogrammierung
wurden im Wintersemester 15/16 das hier beschriebene Betriebssystem f�r Microcontroller entwickelt.
Au�er den funktionalen Vorgaben gab es keine weiteren Bedingungen
oder vordefinierte Codest�cke.
F�r das minimalistische Betriebssystem wurden Module wie der
Dispatcher, Scheduler und Synchronisierungsmechanismen unter
Betracht gezogen. Zudem sollten zwei zus�tliche, frei w�hlbare 
Features implementiert werden. Am Ende stand ein agiles, 
multitasking-f�higes Betriebssystem.
\tableofcontents			% Inhaltsverzeichnis einf�gen
\cleardoublepage

% ==========================================================================
% Inhalt (Einbinden der Dokumente in Ordner Inhalt)
% ==========================================================================
\pagenumbering{arabic} % Seitennummer in arabische Ziffern

\chapter{System�bersicht}

Blabla

\section{Hardware}

\section{Dispatcher}

\section{Scheduler}

\section{Synchronisation}

\section{Zeiten}

\section{Collections}
\chapter{Features}

\section{Speicherverwaltung}
Um dynamische Speicherallokierung zu erm�glichen, was z.B. 
bei der LinkedList n�tig ist, wurde eine Speicherverwaltung implementiert.
Die Speicherverwaltung erlaubt es Speicherbereiche beliebiger Gr��e zu allokieren.
Der von der Speicherverwaltung genutzte Heap wird durch das anlegen eines Segment-Arrays reserviert.
F�r die Segmente ist eine Gr��e von 4 Byte gew�hlt worden.

\subsection{MemoryHead}
Zur Verwaltung der auf dem Heap vergebenen Speicherbereiche wird das vor dem jeweiligen Speicherbereich liegende Segment genutzt.
Dieses Segment wird in einem solchen Fall MemoryHead genannt und ist folgenderma�en aufgebaut (siehe \ref{lst:memoryHead}).

\begin{lstlisting}[language=C, label=lst:memoryHead, caption=MemoryHead Struct]
typedef struct{
	uint16_t prev: 12;
	uint16_t next: 12;
	uint16_t size: 8;
}MemoryHead;
\end{lstlisting}

Die Bitfelder \texttt{prev} und \texttt{next} geben den Index des vorherigen bzw. n�chsten MemoryHead im Heap Array an.
Durch die Limitierung auf 12 Bit k�nnen maximal $2^{12}$ Segmente bzw $4\cdot2^{12} = 16\text{kB}$ addresiert werden.
Eine Implentierung ohne ein \texttt{prev} w�re auch denkbar, also eine einfach verkettete Liste. Es vereinfacht aber das entfernen eines MemoryHeads aus der List wenn ein Speicherbereich freigegeben wird.
Das \texttt{size} Feld gibt die Gr��e des zu verwaltenden Speicherbereichs an.
Da f�r \texttt{size} 8 Bit bereitstehen k�nnen Speicherbereiche mit einer Gr��e von bis zu $4\cdot2^8 = 1\text{kB}$ allokiert werden.

Es existiert zu jeder Zeit mindestens ein MemoryHead: das erste Segment des Heaps.
Dieser dient auch als Einstiegspunkt f�r den Speicherverwaltungsalgorithmus.

\subsection{Algorithmus}
Die Speicherverwaltung legt keinen bestimmten Algorithmus fest, sonder erlaubt es beliebige Algorithmen zu implementieren.
Dazu stellt die Speicherverwaltung zwei Funktionen und eine Hilfsstruktur zur Verf�gung.
Mit der Funktion \texttt{memoryManagement\_next} k�nnen iterativ alle freien Speicherbereiche erfragt werden, die eine angegebene Gr��e nicht unterschreiten.
Die erw�nschte Gr��e wird als Teil der Struktur \texttt{MemoryRequest} �bergeben.
In dieser Struktur werden von der Funktion Informationen �ber den n�chsten freien Speicherbereich gespeichert.
Darunter die Id und die Gr��e des freien Speicherbereichs.
Liefert die Funktion den Wert 0 zur�ck sind keine weiteren freien Speicherbereich der angegebenen Gr��e auf dem Heap verf�gbar.

Um einen Speicherbereich zu belegen steht die Funktion \texttt{memoryManagement\_alloc} zur Verf�gung.
Ihr wird beim Aufrufen ebenfalls ein \texttt{MemoryRequest} �bergeben.
Die in der \texttt{MemoryRequest} Stuktur enthaltene Id gibt dabei an welcher Speicherbereich belegt werden soll.
Der R�ckgabewert der Funktion ist die Adresse des allokierten Speicherbereichs.

Im nachfolgenden Codebeispiel \ref{lst:memAlgoBeispiel} ist der Worst Fit Algorithmus mithilfe der beiden vorgestellten Funktionen implementiert.

\lstinputlisting[language=C, label=lst:memAlgoBeispiel, caption=Worst Fit Algorithmus]{./Quellcode/memAlgoExample.c}

\section{Scheduler: Priorit�tenvererbung}
\label{sect:pip}
Bei der Umsetzung des Priority-Inheritance-Protocols (PIP) wurde 
ein flexibler, selbst-lernender Ansatz erarbeitet. Dabei erweitert
er im Grunde den priorit�tengesteuerten Round-Robin Algorithmus
um zwei Punkte:

\begin{enumerate}
\item Jede Ressource (darunter fallen auch Teile des 
Synchronisations-Moduls) erh�lt jeweils eine eigene Liste.
Darin wird vermerkt, welche Tasks die Ressource jemals freigegeben
haben.
\item Wird nun eine Task durch eine Ressourcenanforderung 
blockiert, �berpr�ft der Algorithmus, ob eine Task aus der
in 1. genannten Liste f�r eine Priorit�tenvererbung in Frage
kommt.
\end{enumerate}

Diese Methode f�hrt dazu, dass nachdem alle Tasks die ben�tigten
Ressourcen mindestens einmal verwendet haben es zu keiner
Priorit�teninversion kommen kann. Soll diese Bedingung gleich
zu Beginn gelten, so kann optional die Freigabeliste bereits
vom Entwickler vorgegeben werden.

\subsection{Blockierung durch Ressourcenanforderung}
Im Fall einer Blockierung wird �ber die o.g. Freigabeliste 
iteriert. Als Bedingung, dass eine Priorit�tenvererbung stattfindet
muss die bevorzugte Task eine niedrigere Priorit�t haben und
lauff�hig sein (\texttt{state = READY}). Ist beides erf�llt,
wird der Task tempor�r die Priorit�t der blockierten Task 
zugewiesen und in die Menge der lauff�higen Tasks neu einsortiert.

\begin{lstlisting}[language=C, label=lst:blockedTask, caption=Verwaltung einer blockierten Ressourcenanforderung]
void prioInheritance_blockedByRessourceRequest(LinkedList* resFreedBy){
	uint8_t length = linkedList_length(resFreedBy);
	taskControlBlock* nextTask;
	if(length){
		for(uint16_t i = 0; i < length; i++){
			linkedList_get(resFreedBy, i, &nextTask);
			if(nextTask->state == READY && nextTask->prio > currentTask->prio){
				nextTask->prio = currentTask->prio;
				queue_removeItem(&prioQueue[nextTask->tmpPrio], nextTask);
				queue_push(&prioQueue[nextTask->prio], nextTask);
				break;
			}
		}
	}
}
\end{lstlisting}

\subsection{Freigabe einer Ressource}
Wird bei PIP eine Ressource freigegeben, wird zun�chst �berpr�ft,
ob die freigebende Task bereits im der Freigabeliste enthalten ist.
Ist dies nicht der Fall, wird sie eingetragen. Wichtiger 
jedoch muss erkannt werden, ob eine Priorit�tenvererbung 
stattgefunden hat. Falls ja, wird die Priorit�t auf den 
urspr�nglichen Wert (\texttt{tmpPrio}) zur�ckgesetzt und die Task
verdr�ngt, da sie nun die Ressource nicht mehr f�r sich in
Anspruch nimmt (s. Code \ref{lst:freedTask} Zeile 7).

\begin{lstlisting}[language=C, label=lst:freedTask, caption=Ressource wird freigegeben]
void prioInheritance_ressourceReleased(LinkedList* resFreedBy){
	if(!linkedList_contains(resFreedBy, currentTask)){
		linkedList_append(resFreedBy, currentTask);
	}
	if(currentTask->tmpPrio != currentTask->prio){
		currentTask->prio = currentTask->tmpPrio;
		task_yield();
	}
}
\end{lstlisting}

\chapter{Tests}
\label{chap:tests}
Um die einzelnen Module des Betriebssystems unabh�ngig testen
zu k�nnen wurde f�r den AVR-Simulator eine eigene Testbench
entwickelt. Mit ihr ist es m�glich s�mtliche Szenarien 
durchzutesten und auf Richtigkeit zu �berpr�fen.

\section{Neustart}
F�r jeden Test sollten dabei eine definierte Umgebung vorliegen,
welche unabh�ngig von zuvor ausgef�hrten Tests sein muss.
Dadurch kann eine flexible Wiederholbarkeit erreicht werden.
Damit alle Tests automatisch ausgef�hrt werden k�nnen wird
ein Betriebssystemneustart ben�tigt. Dazu wird eine Interruptsperre
verh�ngt, der Dispatcher-Timer gestoppt, der Stack geleert, alle
CPU-Register zur�ckgesetzt und schlie�lich zur Programm-Adresse
0 gesprungen.

\begin{lstlisting}[language=C, label=lst:reboot, caption=Neustart]
void beerOS_reboot(){
	disableInterrupts();
	stopDispatcherTimer();
	//clear main stack
	while(mainSP <= RAMEND){
		*mainSP = 0;
		mainSP++;
	}
	mainCalled--;
	SP = 0;
	asm volatile(
		"CLR R0\n\t"
		[...] //clear registers
		"BCLR 7\n\t"
		);
	asm volatile ("jmp 0");
}
\end{lstlisting}
\clearpage

\section{\texttt{.noinit}-Variablen}
Zum Schutz vor einem fehlerhaften mehrfachen Aufrufen der 
\texttt{main}-Funktion wurde ein Z�hler implementiert. Da bei einem
Neustart ein solcher Aufruf gewollt ist wird die Z�hlervariable
\texttt{mainCalled} dekrementiert. Vor dem Start der gewrappten
Startfunktion \texttt{run} wird abgefragt, ob der \texttt{main}-Aufruf
gewollt ist.
\begin{lstlisting}[language=C, label=lst:main, caption=\texttt{main}-Funktion]
int main(void){
	if(mainCalled){
		kernelPanic();
	}
	mainCalled++;
	return run();
}
\end{lstlisting}

Globale Variablen werden vom Compiler in der Weise behandelt, dass
sie beim Systemstart (Durchlauf der Programmadresse 0) 
initialisiert werden. Jedoch ist dieser Automatismus f�r zwei F�lle
nicht gew�nscht:
\begin{enumerate}
\item W�rde die o.g. Variable \texttt{mainCalled} bei jedem Aufruf
der Programmadresse 0 initialisiert bzw. zur�ck gesetzt werden, so w�re die Bedingung in Zeile 2 (Code \ref{lst:main}) nie g�ltig, obwohl
ein Fehlerfall vorl�ge.
\item Auch die Funktionspointer-Variable 
\texttt{void (*initNextTest)(void)}, in der der n�chste 
auszuf�hrende Test gespeichert wird, w�rde zur�ckgesetzt werden.
Automatisiertes Testen w�re somit unm�glich.
\end{enumerate}

Abhilfe hier schafft das Attribut \texttt{(section (".noinit"))}
welches dem Compiler vermittelt, dass diese Variablen zwar 
angelegt, jedoch nicht initialisiert werden sollen.
\footnote{In gewisser Weise werden sie beim ersten Start des 
AVR-Simulators auf die angegebenen Werte initialisiert. Daher 
entsteht beim Compilieren eine Warnmeldung.}

\begin{lstlisting}[language=C, label=lst:noinit, caption=\texttt{noinit}-Variablen]
void (*initNextTest)(void) __attribute__ ((section (".noinit"))) = &initSemaphoreTest;
uint8_t mainCalled __attribute__ ((section (".noinit"))) = 0;
\end{lstlisting}


% ==========================================================================
% Anhang
% ==========================================================================
% \appendix
% \chapter{Erste �berschrift im Anhang}
\label{kap:Anhang:Kapitel}

\section{Unter�berschrift im Anhang}
\label{kap:Anhang:Kapitel:Section}
Lorem ipsum dolor sit amet, consetetur sadipscing elitr, sed diam nonumy eirmod tempor invidunt ut labore et dolore magna aliquyam erat, sed diam voluptua. At vero eos et accusam et justo duo dolores et ea rebum. Stet clita kasd gubergren, no sea takimata sanctus est Lorem ipsum dolor sit amet. Lorem ipsum dolor sit amet, consetetur sadipscing elitr, sed diam nonumy eirmod tempor invidunt ut labore et dolore magna aliquyam erat, sed diam voluptua. At vero eos et accusam et justo duo dolores et ea rebum. Stet clita kasd gubergren, no sea takimata sanctus est Lorem ipsum dolor sit amet.

\chapter{Zweite �berschrift im Anhang}
\label{kap:Anhang:Kapitel2}
Text Text Text Text Text Text Text Text Text Text Text Text Text Text Text Text Text Text Text Text Text Text Text Text Text Text Text Text Text Text Text Text Text Text Text Text Text Text Text Text Text Text Text Text Text Text Text Text Text Text Text Text Text Text Text Text Text Text Text Text Text Text Text Text Text Text Text Text Text Text Text Text Text Text Text Text Text Text Text Text Text Text Text Text Text Text Text Text Text Text.

\section{Nochmal eine Unter�berschrift im Anhang}
\label{kap:Anhang:Kapitel2:Section}
Hier ist ein text

\subsection{eine UnterUnter�berschrift im Anhang}
\label{kap:Anhang:Kapitel2:Section:Subsection}
Hier ist ein text

% \cleardoublepage

% ==========================================================================
%	Literatur- Abbildung- und Abk�rzungsverzeichnis
% ==========================================================================
%Abk�rzungen ausgeben
\printglossary[type=\acronymtype,title=Abk�rzungsverzeichnis,toctitle=Abk�rzungsverzeichnis,style=myacro]

%Symbole ausgeben
\printglossary[type=symbolslist,title=Symbolverzeichnis,toctitle=Symbolverzeichnis,style=mysombol]

\listoftables					% Tabellenverzeichnis

\listoffigures				% Abbildungsverzeichnis

% \begin{RaggedRight} %Darstellung des Literaturverzeichnisses im Rausatz wegen URLs
% \renewcommand{\bibname}{Quellenverzeichnis}	% Literaturverzeichnis umbenennen
% \bibliography{Dokumentation}    %Dateiname der .bib Datei
% \end{RaggedRight}
\clearpage

% ==========================================================================
%	ToDo-Liste (kann auskommentiert werden, wenn nicht ben�tigt)
% ==========================================================================
\listoftodos[To-do Liste]
\todototoc

\end{document}
